\documentclass[]{article}

\usepackage{amsmath}
\usepackage{multirow}
\usepackage{natbib}
\usepackage{graphicx} 


\begin{document}

The development of large-scale carbon capture technologies is critical for mitigating the impact of anthropogenic CO₂ emissions from industrial point sources. Solid sorbents offer a promising pathway for capturing CO₂ from flue gas streams, but their practical implementation hinges on satisfying a set of demanding and often conflicting material properties. An ideal sorbent must not only exhibit a strong chemical affinity for CO₂ but also release it with a minimal energy penalty during thermal regeneration. Furthermore, to be economically viable, the material must maintain its structural and chemical integrity over thousands of high-temperature capture-and-release cycles.

The central challenge in sorbent design lies in navigating the inherent trade-offs between these requirements. The thermodynamics of the carbonation reaction must be precisely tuned: materials that bind CO₂ too strongly, such as many alkali and alkaline earth oxides, require excessively high regeneration temperatures, rendering the process energetically inefficient. Conversely, materials with weak CO₂ binding may be unreactive under typical flue gas conditions. Beyond this thermodynamic dilemma, sorbents face severe degradation mechanisms. The solid-state transformation from an oxide to a carbonate is often accompanied by significant volume changes, which can induce particle fracture, loss of active surface area, and a rapid decline in performance. Simultaneously, the elevated temperatures needed for regeneration can cause sintering, where particles fuse and further diminish the material's capture capacity. The search for a single material that simultaneously resolves these competing thermodynamic, mechanical, and thermal challenges has remained a primary bottleneck.

To accelerate the discovery of durable sorbents, we present a high-throughput computational screening strategy that systematically evaluates a vast chemical space. Leveraging the Materials Project database, we computationally construct and analyze 889 unique metal oxide-carbonate reaction cycles. Our approach begins by systematically parsing chemical formulas to identify all stoichiometrically balanced reaction pairs, ensuring a consistent basis for thermodynamic comparison. We then filter this dataset to include only materials that are either on the convex hull of thermodynamic stability or close to it, focusing the search on experimentally accessible compounds. Each candidate is subsequently assessed against a multi-faceted set of performance metrics designed to address the key failure modes. The Gibbs free energy of the carbonation reaction, $\Delta G$, is calculated to identify materials that operate within the optimal thermodynamic window for reversible cycling. The volumetric expansion upon carbonation is quantified to filter out materials susceptible to mechanical breakdown. Finally, the Tamman temperature is used as a proxy to estimate the material's resistance to sintering at operational temperatures.

This comprehensive screening reveals that simple binary oxides are generally poor candidates, as they tend to occupy thermodynamic extremes of being either too stable for regeneration or too unreactive for capture. We find that catastrophic volumetric expansion is a dominant and often overlooked failure mechanism, disqualifying a large fraction of otherwise promising materials. The materials that successfully emerge from our rigorous filtering process are not simple compounds but are overwhelmingly complex, mixed-metal polyanionic frameworks. By balancing these competing requirements, candidates such as sodium titanium phosphates and lithium vanadium phosphates are identified as a new class of durable sorbents that exhibit moderate reaction thermodynamics, minimal volume change, and high predicted thermal stability. This work provides a systematic framework for materials discovery and pinpoints a promising chemical space for the rational design of robust and efficient CO₂ capture technologies.

\end{document}
                